\section{Related Work}\label{sec:related}

\subsection{Genotype, subtype, and morphology}\label{sec:related/morphology}

The literature supports a hierarchical view of the genotype-morphology relationship, from near-obligate, mechanistically-traceable lesions at the top to point mutations with no recognizable imprint at the bottom.

At the top, a small set of genes produce near-pathognomonic morphology in their tissue context.
CDH1 loss disrupts adherens junctions, producing single-cell discohesion characteristic of invasive lobular carcinoma~\cite{frixenEcadherinmediatedCellcellAdhesion1991, derksenSomaticInactivationEcadherin2006}.
VHL loss stabilizes HIF, which induces PLIN2 to package lipid droplets~\cite{qiuHIF2aDependentLipidStorage2015} and represses CPT1A to block fatty acid oxidation~\cite{duHIFDrivesLipid2017}; together with HIF-driven glycogen accumulation~\cite{chenGYS1InducesGlycogen2020, pescadorHypoxiaPromotesGlycogen2010}, this produces the clear cytoplasm of clear cell renal cell carcinoma.
PML-RAR$\alpha$ arrests myeloid maturation at the promyelocyte stage through enhanced corepressor/HDAC recruitment~\cite{grignaniFusionProteinsRetinoic1998, linRoleHistoneDeacetylase1998}, producing the hypergranular promyelocytes with faggot cells characteristic of acute promyelocytic leukemia~\cite{swerdlowWHOClassificationTumours, sehgalAuerRodsFaggot2023}.
MYC deregulation couples proliferation with apoptosis~\cite{evanInductionApoptosisFibroblasts1992, eischenDisruptionARFMdm21999}; combined with cooperating mutations, it produces Burkitt lymphoma~\cite{schmitzBurkittLymphomaPathogenesis2012}, where high cell turnover generates the characteristic starry-sky pattern of tingible-body macrophages clearing apoptotic cells~\cite{swerdlowWHOClassificationTumours, dy-ledesmaStarrySkyPattern2016}.
The lesion–morphology association recurs across many tumor types:
the DNAJB1-PRKACA fusion is present in essentially all fibrolamellar hepatocellular carcinomas~\cite{honeymanDetectionRecurrentDNAJB1PRKACA2014}, which are defined by their lamellar-fibrosis architecture~\cite{craigFibrolamellarCarcinomaLiver1980}; 
germline FH mutation followed by somatic loss of the second allele~\cite{tomlinsonGermlineMutationsFH2002} drives fumarate accumulation~\cite{pollardAccumulationKrebsCycle2005, adamRenalCystFormation2011} and is associated with the inclusion-like nucleolar morphology recognized as the hallmark of HLRCC RCC~\cite{merinoMorphologicSpectrumKidney2007}; 
the ETV6-NTRK3 fusion produces the microcystic, secretion-rich morphology shared by secretory breast carcinoma~\cite{tognonExpressionETV6NTRK3Gene2002} and its salivary-gland counterpart, mammary analogue secretory carcinoma (MASC)~\cite{skalovaMammaryAnalogueSecretory2010}.
In each case a single defining lesion, which either disrupts a tissue-specific differentiation program or imposes an obligatory metabolic or architectural phenotype, is associated with morphology that a pathologist can recognize at low magnification.

Beneath this layer, the link weakens from obligate to probabilistic: a subtype enriches for a recognizable morphologic cluster, but the same morphology can arise outside it and a member of that subtype can present without it. 
In breast cancer, the PAM50 intrinsic subtypes~\cite{parkerSupervisedRiskPredictor2009} partition tumors along an axis of luminal differentiation: luminal tumors retain ESR1, GATA3, and FOXA1 expression, while basal-like tumors lose this program and present as high-grade carcinomas~\cite{koboldtComprehensiveMolecularPortraits2012} with medullary-like features~\cite{livasyPhenotypicEvaluationBasallike2006}.
Deep learning can partially recover this signal from H\&E at $\approx$77\% accuracy for the binary basal-vs-non-basal call~\cite{coutureImageAnalysisDeep2018} and $\approx$66\% for the four-class call~\cite{jaberDeepLearningImagebased2020}.
In colorectal cancer, the four consensus molecular subtypes, including MSI immune (CMS1), WNT/MYC canonical (CMS2), KRAS metabolic (CMS3), and TGF-$\beta$ mesenchymal (CMS4)~\cite{guinneyConsensusMolecularSubtypes2015}, leave H\&E imprints (Crohn-like lymphoid infiltrates in CMS1~\cite{greensonPathologicPredictorsMicrosatellite2009}; desmoplastic stroma in CMS4~\cite{isellaStromalContributionColorectal2015}) that the imCMS classifier~\cite{sirinukunwattanaImagebasedConsensusMolecular2021} predicts at macro-AUROC $\approx$0.84 on held-out cohorts.
In endometrial cancer, the four TCGA molecular classes cluster along axes of mutational load and chromosomal stability: POLE-ultramutated and MMRd hypermutated tumors carry high neoantigen loads and dense TILs~\cite{howittAssociationPolymeraseEMutated2015}, while p53-abnormal copy-number-high tumors display marked nuclear atypia and serous-like architecture~\cite{levineIntegratedGenomicCharacterization2013, fremondInterpretableDeepLearning2023}. Pathologist discordance motivated the molecular reclassification~\cite{gilksPoorInterobserverReproducibility2013}; all four classes can now be predicted from H\&E~\cite{fremondInterpretableDeepLearning2023} as a component of the HECTOR recurrence prognosticator~\cite{volinsky-fremondPredictionRecurrenceRisk2024}.
The pattern recurs in thyroid cancer, where BRAF V600E drives strong MAPK activation and the classical-type papillary architecture while RAS-family mutations produce weaker MAPK signaling and a follicular-pattern morphology~\cite{agrawalIntegratedGenomicCharacterization2014}, and in small cell lung cancer, where the ASCL1-, NEUROD1-, POU2F3-, and YAP1-defined transcriptional subtypes are defined transcriptionally rather than morphologically~\cite{rudinMolecularSubtypesSmall2019}.
In each case the morphology emerges from the aggregate of co-occurring alterations and the transcriptional program they engage; no single mutation suffices to produce it.
This is why subtype prediction from H\&E achieves consistently higher within-cohort AUROCs than individual-mutation prediction, even when the same model architecture is used for both targets~\cite{katherPancancerImagebasedDetection2020, arslanSystematicPancancerStudy2024}.

At the bottom of the hierarchy, most somatic point mutations have no human-recognizable morphologic signature and are predictable from H\&E only via cell-of-origin, subtype enrichment, or coincidental association with a more morphologically penetrant alteration.
Pan-cancer integrative analyses confirm that cell-of-origin, not individual driver mutations, dominates the molecular taxonomy of human tumors~\cite{hoadleyCellofOriginPatternsDominate2018}, and the same asymmetry is observed when the predictor is H\&E: pan-cancer studies that train one model per recurrent driver mutation report AUROCs clustered in the 0.6--0.8 range, with mean single-nucleotide-variant performance materially lower than subtype prediction using the same architectures~\cite{coudrayClassificationMutationPrediction2018, fuPancancerComputationalHistopathology2020, katherPancancerImagebasedDetection2020, arslanSystematicPancancerStudy2024}.
Where mutation prediction does work, the signal commonly reflects a confounder rather than mutation-specific morphology: TP53 status in breast cancer tracks the basal-like subtype it is enriched in, so apparent ``TP53 prediction'' is largely subtype prediction~\cite{koboldtComprehensiveMolecularPortraits2012}; in colorectal cancer, the same pipeline that generalizes well for MSI and BRAF fails for KRAS, NRAS, and PIK3CA on external validation~\cite{niehuesGeneralizableBiomarkerPrediction2023}; and gene-level predictors can collapse onto site- or stain-specific shortcuts not visible to a pathologist~\cite{niehuesGeneralizableBiomarkerPrediction2023, howardImpactSitespecificDigital2021}.
The apparent exceptions are precisely the cases where the mutation drives a broader subtype state with its own morphology, such as MSI in colorectal cancer (functionally CMS1)~\cite{echleClinicalGradeDetectionMicrosatellite2020,guinneyConsensusMolecularSubtypes2015}, BRAF V600E in thyroid (driving classical papillary architecture)~\cite{agrawalIntegratedGenomicCharacterization2014}, so the ``mutation predictor'' is operating at the subtype level.
A consequence, central to the present work, is that any single somatic point mutation is unlikely to leave a reproducible H\&E imprint on its own.

A consequence of this hierarchy, central to the present work, is that any single somatic point mutation is unlikely to leave a reproducible H\&E imprint on its own. 
This motivates predicting morphology from \emph{summary} representations of the somatic genome, such as subtype assignments and mutational signature exposures, rather than from the presence or absence of individual point mutations.
Subtype prediction from H\&E is already a mature direction~\cite{coutureImageAnalysisDeep2018, jaberDeepLearningImagebased2020, sirinukunwattanaImagebasedConsensusMolecular2021, fremondInterpretableDeepLearning2023, volinsky-fremondPredictionRecurrenceRisk2024}; 
signature-exposure prediction is younger but the most clinically actionable summary, homologous-recombination deficiency, can be predicted from H\&E with strong, bias-controlled discrimination (AUROC 0.83--0.86 after correcting for technical batch and molecular-subtype confounders) in luminal breast cancer~\cite{lazardDeepLearningIdentifies2022}.

\subsection{Deep learning on whole slide images}\label{sec:related/dl-wsi}

\autoref{sec:intro} traced the broader trajectory of deep learning on H\&E images, from the CNN-MIL era to large self-supervised foundation encoders.
This subsection focuses on the methodological lineage most directly relevant to Augur: weakly-supervised aggregation of tile features into slide-level predictions, multi-task and pretext supervision for tile feature extraction, and attention-based explainability.

\paragraph{Weakly-supervised aggregation}
Whole-slide images are gigapixel objects, routinely $100{,}000 \times 100{,}000$ pixels or more. They cannot be processed by a CNN in a single forward pass, while slide-level diagnoses are abundant in pathology reports and pixel-level annotations are not.
The standard response partitions each slide into thousands of tiles, encodes them independently into feature vectors, and aggregates those features into a single slide-level prediction. The formal substrate is multiple-instance learning (MIL), which treats a slide as a bag of instances carrying a single bag-level label and parameterizes the Bernoulli bag-label distribution end-to-end via a tile encoder, a permutation-invariant pooling operator, and a bag classifier~\cite{ilseAttentionbasedDeepMultiple2018}.
Campanella et al.~\cite{campanellaClinicalgradeComputationalPathology2019} trained max-pool MIL models on 44{,}732 H\&E WSIs from 15{,}187 patients across three cancer types using only report-derived slide-level labels, reaching test AUROCs of 0.97--0.99 across the three tasks and outperforming pixel-annotated fully-supervised baselines on cross-cohort transfer.

The choice of pooling operator is consequential: fixed max- or mean-pooling collapses the bag to a single score without indicating which tiles contributed, sacrificing both flexibility and interpretability. Attention-based MIL replaces fixed pooling with a learned attention distribution over tiles, yielding both higher accuracy in the small-sample regime and per-tile saliency maps that align with ground-truth annotations under bag-only supervision~\cite{ilseAttentionbasedDeepMultiple2018}. 
CLAM~\cite{luDataefficientWeaklySupervised2021} extends this with an instance-clustering auxiliary loss that pushes high-attention tiles to be class-discriminative and low-attention tiles to be uninformative; its CLAM-MB variant additionally allocates one attention head per output class for class-conditional saliency, and CLAM reaches 10-fold macro-AUROCs of 0.991 / 0.956 / 0.953 on renal cell carcinoma subtyping, NSCLC subtyping, and lymph-node metastasis detection, while retaining AUROC $>$0.9 with as little as 25--50\% of the training data.
TransMIL~\cite{shaoTransMILTransformerBased2021} replaces the i.i.d.-tile assumption with correlated MIL realized through self-attention with multi-scale positional encoding, achieving AUROCs of 0.93 / 0.96 / 0.99 on CAMELYON16 / TCGA-NSCLC / TCGA-RCC, roughly 5 AUROC points above CLAM on CAMELYON16, with the trade-off that the tile-to-tile attention matrix is harder to read as per-tile saliency than the scalar attention weights of gated MIL.

\paragraph{Multi-task and pretext supervision}

Early supervised pathology pipelines trained one model per task (tile classification, nuclear segmentation, mutation prediction, etc.), each requiring dense human annotation that scaled poorly.
Pan-cancer mutation predictors require matched H\&E and sequencing data for every recurrent driver and reach only modest AUROCs (0.73--0.86 across NSCLC drivers; only 6 of 10 tested genes reached significance)~\cite{coudrayClassificationMutationPrediction2018}; nuclear-instance models require exhaustive per-instance labels in the thousands.

Multi-task supervision amortizes annotation cost across tasks that share a feature substrate. Within a single image, HoVer-Net~\cite{grahamHoverNetSimultaneousSegmentation2019} allocates three decoder branches to nuclear pixel segmentation, horizontal-vertical distance regression, and per-nucleus classification, producing simultaneous state-of-the-art performance on each (CoNSeP PQ 0.547, classification F1 0.42--0.64 across nuclear types) on a dataset of 24{,}319 manually annotated nuclei.
Its Cerberus~\cite{grahamOneModelAll2023} extension scales the same principle to nuclei, glands, lumina, and tissue regions, exceeding HoVer-Net on every shared task (internal PQ 0.612 vs 0.583, external mPQ 0.332 vs 0.285) and showing that the shared multi-task encoder transfers to downstream tasks not seen during training.
At the cross-dataset level, Mormont et al.~\cite{mormontMultiTaskPreTrainingDeep2021} pretrain a single backbone jointly on 22 histopathology classification tasks ($\sim$900k images), beating ImageNet on 15 of 20 downstream evaluations under frozen feature extraction, although the advantage disappears under full fine-tuning, which lets ImageNet's generic features catch up.

Self-supervised learning takes the annotation burden off entirely by deriving supervision from the input itself. 
In Self-Path framework~\cite{koohbananiSelfPathSelfsupervisionClassification2020}, domain-specific pretexts tailored to pathology (magnification prediction, magnification-jigsaw task (JigMag), and hematoxylin-channel reconstruction) outperform supervised baselines by 6--13 AUC points with as little as 1\% of WSI labels.
Generic contrastive learning (SimCLR) applied to a 57-dataset histopathology corpus boosts downstream linear-probe macro-F1 by $\sim$28 points over ImageNet pretraining, with diminishing returns past $\sim$50k unlabelled images~\cite{cigaSelfSupervisedContrastive2022}.
Transformer-based SSL with semantically-mined positives (CTransPath/SRCL) pretrains a hybrid CNN-Swin backbone on $\sim$15.6M TCGA and PAIP patches and, paired with CLAM, reaches WSI-level AUROCs of 0.922 on Camelyon16, 0.912 on TCGA-NSCLC, and 0.967 on TCGA-RCC~\cite{wangTransformerbasedUnsupervisedContrastive2022}.
Masked image modeling at TCGA scale (Phikon, iBOT ViT-B on the PanCancer40M corpus of 43M tiles) extends this further, beating CTransPath by +1.4\% and HIPT by +4.0\% on average across 14 weakly-supervised WSI tasks and reaching 98.2 AUROC on PAIP-CRC MSI external validation~\cite{filiotScalingSelfSupervisedLearning2024}.

\paragraph{Explainability}



% \subsection{Mutational landmarks in breast cancer}
% \label{sec:related/breast-mutations}

% Breast cancer is the most extensively studied indication for H\&E-based molecular prediction, and provides the closest precedent for the present work.
% We organize prior efforts by the molecular target: intrinsic subtypes, individual driver-gene status, and aggregate genomic-instability features --- with the last category most directly overlapping VexDR.

% \paragraph{Intrinsic subtypes}
% Couture et al.~\cite{coutureImageAnalysisDeep2018} were among the first to predict the four-class PAM50 intrinsic subtype from H\&E using CNNs trained on tissue microarrays, reporting basal-versus-non-basal accuracy near 77\% and ductal-versus-lobular accuracy near 94\%.
% Naik et al.~\cite{naikDeepLearningenabledBreast2020} extended this to whole-slide images and showed that within-slide heterogeneity in predicted PAM50 subtype is itself prognostically informative.
% Subsequent foundation-model-based subtypers~\cite{luVisuallanguageFoundationModel2024, chenGeneralpurposeFoundationModel2024} have pushed subtype AUROCs above 0.9 in zero-shot or few-shot settings.
% The asymmetry highlighted in Section~\ref{sec:intro} --- that subtypes are markedly more predictable than any single mutation they enrich for --- is now an empirical regularity across breast cancer studies and underlies our choice of subtype as the \emph{main} task in VexDR rather than the pretext.

% \paragraph{Individual driver-gene status}
% Pan-cancer mutation predictors~\cite{coudrayClassificationMutationPrediction2018, katherPancancerImagebasedDetection2020, fuPancancerComputationalHistopathology2020} typically report breast-cancer AUROCs in the 0.6--0.8 range for individual recurrent drivers (TP53, PIK3CA, GATA3).
% TP53 has received specific attention because TP53-mutant breast cancers are enriched for high grade, basal-like phenotype, and chromosomal instability: dedicated TP53 predictors trained on TCGA-BRCA achieve moderate AUROCs but a recurring concern is that the predicted signal partly reflects the basal-like context in which TP53 mutations are enriched, rather than a TP53-specific morphology~\cite{naikDeepLearningenabledBreast2020}.
% Within the genotype-morphology hierarchy of Section~\ref{sec:related/genotype-morphology}, this observation places most individual point-mutation predictors at the bottom layer: they recover predominantly cell-of-origin and subtype context rather than mutation-specific morphology.

% \paragraph{Aggregate genomic-instability features and mutational signatures}
% The closest precedent for VexDR is the line of work that predicts \emph{aggregate} mutational-process features from H\&E, most prominently homologous-recombination deficiency (HRD).
% HRD is itself a composite phenotype: it is most often defined operationally by the activity of COSMIC SBS3 (and the indel signature ID6), by genomic scar metrics (LOH, telomeric allelic imbalance, large-scale state transitions), or by combined HRD scores such as myChoice and Genomic LOH.
% Lazard et al.~\cite{lazardDeepLearningIdentifies2022} train a deep MIL model to predict HRD status from H\&E in luminal breast cancers, after explicitly controlling for confounding by grade and tumor cellularity, and identify laminated fibrosis and clear tumor cell foci as morphological correlates of HRD.
% DeepHRD (Bergstrom et al.~\cite{bergstromDeepLearningArtificial2024}) is a weakly-supervised MIL HRD predictor trained on TCGA-BRCA H\&E with external validation across breast and ovarian cohorts, reaching AUROC 0.81 and outperforming FDA-approved molecular HRD assays for predicting platinum-specific PFS.
% The SuRe-Transformer~\cite{huBreastCancerHomologous2025} pushes HRD AUROC to 0.887 with a transformer-MIL backbone, and Schirris et al.~\cite{schirrisPredictionHomologousRecombination2024} demonstrate that an attention-MIL HRD predictor generalizes across nine tumor types.
% Wagner et al.~\cite{wagnerTransformerbasedBiomarkerPrediction2023} show that regression-based MIL is preferable to thresholded classification for HRD-like continuous targets.

% A second well-studied mutational process is APOBEC mutagenesis (COSMIC SBS2/SBS13), which is enriched in HER2-amplified breast cancers and in invasive lobular carcinoma~\cite{robertsAPOBECCytidineDeaminase2013}, suggesting that APOBEC activity has at least an indirect morphologic correlate via its subtype enrichment.
% A small number of studies have begun to predict APOBEC-related summary scores from H\&E, but these have generally treated APOBEC as a binary high/low call rather than as a component of a multi-process exposure profile.
% Microsatellite instability and POL$\epsilon$-driven hypermutation, both of which generate distinctive SBS signatures (SBS6/15/20/26 for MMR-deficiency; SBS10a/b for POL$\epsilon$), have been predicted from H\&E primarily in colorectal and endometrial cancers; their breast-cancer prevalence is too low for direct training on TCGA-BRCA.

% \paragraph{Where VexDR sits}
% Three observations follow from the literature reviewed above.
% First, prior work has shown convincingly that \emph{individual} mutational-process summaries --- HRD status, MSI status, TMB, occasionally APOBEC --- can be predicted from H\&E with clinically useful discrimination, validating the broader premise that summary genomic features leave reproducible morphologic imprints.
% Second, this prior work has overwhelmingly treated each summary as a stand-alone, often-binarized target; we are not aware of any breast-cancer study that jointly predicts the \emph{full SBS exposure vector} (multiple co-active processes, on a continuous scale) alongside intrinsic subtype, with attention disentangled per process.
% Third, available work treats genotype prediction as the deployment goal; VexDR instead positions SBS exposure prediction as an explanatory pretext on top of the clinically established subtype-classification task, with the explicit aim of using process-resolved attention maps as a hypothesis-generation tool for novel morphology-process associations not previously documented in wet-lab studies.

% % ============================================================================
% %                         ARCHIVE TEXT
% % ============================================================================
% Beneath this layer, the link weakens from obligate to probabilistic: a tumor in a given subtype is enriched for a recognizable morphologic cluster, but the same morphology can arise outside the subtype, and a tumor of that subtype can present without it. 
% In breast cancer, the PAM50 intrinsic subtypes~\cite{parkerSupervisedRiskPredictor2009} partition tumors along an axis of luminal differentiation: luminal tumors retain expression of ESR1, GATA3, and FOXA1 and are typically low-to-intermediate grade, while basal-like tumors lose this program (with frequent TP53 inactivation and basal-cytokeratin expression) and present as high-grade carcinomas~\cite{koboldtComprehensiveMolecularPortraits2012}, often with medullary-like features including geographic necrosis and prominent tumor-infiltrating lymphocytes~\cite{livasyPhenotypicEvaluationBasallike2006}. 
% Deep learning models can partially recover this molecular signal from H\&E images: as a binary basal-vs-non-basal call with $\approx$77\% accuracy~\cite{coutureImageAnalysisDeep2018}, and as a four-class call with $\approx$66\% accuracy~\cite{jaberDeepLearningImagebased2020}.
% In colorectal cancer, the four consensus molecular subtypes (CMS1--4)~\cite{guinneyConsensusMolecularSubtypes2015} are characterized by distinct molecular programs: MSI-driven hypermutation and immune activation in CMS1, WNT/MYC signaling in CMS2, KRAS-driven metabolic dysregulation in CMS3, and TGF-$\beta$ activation with stromal infiltration in CMS4. These programs leave morphologic signatures recognizable on H\&E:\@
% CMS1 tumors show dense lymphocytic infiltrates with Crohn-like peritumoral reaction~\cite{greensonPathologicPredictorsMicrosatellite2009}, while CMS4 tumors show prominent desmoplastic stroma and a stromally-driven mesenchymal phenotype~\cite{isellaStromalContributionColorectal2015}.
% The imCMS classifier of Sirinukunwattana et al.~\cite{sirinukunwattanaImagebasedConsensusMolecular2021} makes this morpho-molecular link explicit, predicting CMS labels from H\&E sections at macro-AUROC $\approx$0.84 on held-out cohorts.
% In endometrial cancer, the four TCGA molecular classes cluster along axes of mutational load and chromosomal stability: POLE-ultramutated and MSI/MMRd hypermutated tumors carry high neoantigen loads and dense TIL infiltrates~\cite{howittAssociationPolymeraseEMutated2015}, while p53-abnormal copy-number-high tumors display the marked nuclear atypia and serous-like architecture characteristic of extensive somatic copy-number alteration~\cite{levineIntegratedGenomicCharacterization2013, fremondInterpretableDeepLearning2023}.
% Pathologist discordance in distinguishing high-grade endometrioid from serous histology motivated the molecular reclassification~\cite{gilksPoorInterobserverReproducibility2013}, and recent work has shown that all four classes can be approximated directly from H\&E using deep learning (im4MEC,~\cite{fremondInterpretableDeepLearning2023}), which forms a component of the HECTOR recurrence prognosticator of Volinsky-Fremond et al.~\cite{volinsky-fremondPredictionRecurrenceRisk2024}.
% The same pattern recurs in thyroid cancer, BRAF V600E drives strong MAPK activation and the classical-type papillary architecture of conventional papillary thyroid carcinoma, whereas RAS-family mutations produce weaker MAPK signaling and a predominantly follicular-pattern morphology~\cite{agrawalIntegratedGenomicCharacterization2014}; and in small cell lung cancer, where the ASCL1-, NEUROD1-, POU2F3-, and YAP1-defined transcriptional subtypes reflect distinct cell-of-origin and differentiation programs, though they are defined transcriptionally rather than morphologically~\cite{rudinMolecularSubtypesSmall2019}.
% In each case the morphology emerges from the aggregate effect of co-occurring alterations and the transcriptional program they engage; no single mutation suffices to produce it. 
% This is why subtype prediction from H\&E achieves consistently higher within-cohort AUROCs than individual-mutation prediction, even when the same model architecture is used for both targets~\cite{katherPancancerImagebasedDetection2020, arslanSystematicPancancerStudy2024}.